% ============================================================================
%  preamble/common.tex — style dùng chung cho toàn cây marker.
%  Dẫn xuất nguyên vẹn từ preamble của cây IMU_Fusion (cùng tác giả, cùng quy ước);
%  chỉ đổi bảng repo ở dưới sang hệ sinh thái fiducial marker, và nhãn mini project.
%  Biên dịch: XeLaTeX + biber. Mọi macro, hộp, font định nghĩa ở đây.
%  KHÔNG định nghĩa lại bất cứ thứ gì ở file con.
% ============================================================================
\usepackage[a4paper,top=2.3cm,bottom=2.4cm,left=2.4cm,right=2.4cm,headsep=0.6cm]{geometry}

% ---------- font, ngôn ngữ ----------
\usepackage{fontspec}
\setmainfont{TeX Gyre Pagella}
\setsansfont{Lato}[Scale=MatchLowercase, BoldFont={Lato Semibold}]
\setmonofont{Noto Sans Mono}[Scale=0.86]
\usepackage{amsmath,amssymb,mathtools}
\usepackage[vietnamese]{babel}

% ---------- màu ----------
\usepackage[table]{xcolor}
\definecolor{inkblue}{HTML}{12395B}     % xanh navy sâu — tiêu đề
\definecolor{steel}{HTML}{2E7290}       % xanh teal — tiêu đề phụ, nhãn
\definecolor{accent}{HTML}{C0651C}      % hổ phách — số chương, chữ phần, điểm nhấn
\definecolor{rulecol}{HTML}{CBD7E1}
\definecolor{violetdark}{HTML}{5A3E8C}
\definecolor{violetbg}{HTML}{F4F1FA}
\definecolor{reddark}{HTML}{A63030}
\definecolor{redbg}{HTML}{FDF1F0}
\definecolor{bluebg}{HTML}{EDF4F9}
\definecolor{codebg}{HTML}{F7F7F4}
\definecolor{tiercol}{HTML}{6E7C89}
\definecolor{panelbg}{HTML}{F4F7FA}     % nền khối tóm tắt
\definecolor{amberbg}{HTML}{FDF5EC}     % nền khối kiến thức nền
\definecolor{tblhead}{HTML}{E7EEF4}     % nền hàng tiêu đề bảng
\definecolor{tblalt}{HTML}{F7FAFC}      % sọc bảng
\definecolor{tblrule}{HTML}{8FA6B8}     % đường kẻ bảng
\definecolor{greendark}{HTML}{2F6B4F}   % khối nói lại thật đơn giản
\definecolor{greenbg}{HTML}{EFF6F1}

% ---------- bảng, danh sách, hình ----------
\usepackage{array,longtable,tabularx,booktabs,multirow,makecell}
\usepackage{etoolbox}
% bảng: sọc nhẹ từ hàng thứ hai, đường kẻ màu xanh xám thay vì đen
\arrayrulecolor{tblrule}
\BeforeBeginEnvironment{tabularx}{\rowcolors{2}{white}{tblalt}}
\BeforeBeginEnvironment{tabular}{\rowcolors{2}{white}{tblalt}}
\BeforeBeginEnvironment{longtable}{\rowcolors{2}{white}{tblalt}}
\setlength{\heavyrulewidth}{1.1pt}
\setlength{\lightrulewidth}{0.5pt}
\usepackage{enumitem}
\setlist{leftmargin=1.5em,itemsep=0.22em,topsep=0.4em,parsep=0pt}
\usepackage{microtype}
\usepackage{graphicx}
\usepackage{float}
\usepackage{tikz}
\usetikzlibrary{arrows.meta,positioning,calc,fit,backgrounds,decorations.pathreplacing,shapes.geometric,patterns}
\usepackage{pgfplots}\pgfplotsset{compat=1.18}
\usepackage[most]{tcolorbox}
\usepackage{emptypage}
\usepackage{listings}
\usepackage{colortbl}
\usepackage[font=small,labelfont={bf,sf},labelsep=period,skip=6pt]{caption}
\usepackage{titletoc}
\usepackage[all]{nowidow}

\newcolumntype{L}[1]{>{\raggedright\arraybackslash}p{#1}}
\newcolumntype{Y}{>{\raggedright\arraybackslash}X}
\renewcommand{\arraystretch}{1.18}
\setlength{\tabcolsep}{5pt}

% tiếng Việt không có mẫu ngắt từ; cho TeX chỗ co giãn
\tolerance=3000 \emergencystretch=3.5em \hbadness=7000
\setlength{\parskip}{0.5em}
\setlength{\parindent}{0pt}

% ---------- tiêu đề ----------
\usepackage{titlesec}

% PHẦN: trang phân cách có chữ cái lớn mờ làm nền
\titleformat{\part}[display]
  {\normalfont\sffamily\bfseries\color{inkblue}\raggedright}
  {\hfill\begin{tikzpicture}[remember picture,overlay]
     \node[anchor=north east,inner sep=0pt] at ($(current page.north east)+(-2.2cm,-4.2cm)$)
       {\fontsize{140}{140}\selectfont\sffamily\bfseries\color{accent!16}\thepart};
   \end{tikzpicture}%
   {\large\sffamily\bfseries\color{accent}PHẦN \thepart}}
  {0.4em}{\fontsize{30}{34}\selectfont\raggedright}
  [\vspace{0.7em}{\color{accent}\rule{5cm}{2pt}}]
\titlespacing*{\part}{0pt}{3.2cm}{2.2em}

% CHƯƠNG: số lớn màu hổ phách bên phải, tiêu đề lớn, gạch mảnh
\titleformat{\chapter}[display]
  {\normalfont\sffamily\bfseries\color{inkblue}}
  {\hfill{\fontsize{46}{46}\selectfont\color{accent!75}\thechapter}}
  {-0.15em}
  {\raggedright\fontsize{21}{25}\selectfont}
  [\vspace{0.45em}{\color{rulecol}\titlerule[1.2pt]}]
\titlespacing*{\chapter}{0pt}{0.4em}{1.7em}

% MỤC: thanh nhấn nhỏ trước số mục
\titleformat{\section}
  {\normalfont\sffamily\large\bfseries\color{inkblue}}
  {\raisebox{0.06em}{\color{accent}\rule{2.6pt}{0.82em}}\hspace{0.5em}\thesection}
  {0.55em}{}
  [\vspace{-0.42em}{\color{rulecol}\rule{\linewidth}{0.7pt}}]
\titlespacing*{\section}{0pt}{1.7em}{0.65em}

\titleformat{\subsection}{\normalfont\sffamily\bfseries\color{steel}}{\thesubsection}{0.5em}{}
\titlespacing*{\subsection}{0pt}{1.15em}{0.3em}
\titleformat{\subsubsection}{\normalfont\sffamily\bfseries\itshape\color{steel}}{}{0pt}{}
\titlespacing*{\subsubsection}{0pt}{0.9em}{0.2em}
\setcounter{secnumdepth}{1}   % chỉ đánh số tới \section; subsection là tiêu đề không số
\setcounter{tocdepth}{1}
\renewcommand{\thepart}{\Alph{part}}
% khi compile riêng một file con, chapter=0 nên bỏ tiền tố chương
\renewcommand{\thesection}{\ifnum\value{chapter}>0 \thechapter.\fi\arabic{section}}

% mục lục: giãn cách rõ ràng hơn
\titlecontents{chapter}[2.6em]{\vspace{0.5em}\sffamily\bfseries\color{inkblue}}
  {\contentslabel[\thecontentslabel]{2.4em}}{\hspace*{-2.4em}}
  {\hfill\normalfont\color{tiercol}\contentspage}
\titlecontents{section}[5.0em]{\smallskip\small}
  {\contentslabel[\thecontentslabel]{2.6em}}{}
  {\titlerule*[0.6pc]{\color{rulecol}.}\small\color{tiercol}\contentspage}

\usepackage{fancyhdr}
\pagestyle{fancy}\fancyhf{}
\renewcommand{\headrulewidth}{0pt}
\fancyhead[L]{\footnotesize\sffamily\color{tiercol}\nouppercase{\leftmark}}
\fancyhead[R]{\small\sffamily\bfseries\color{white}%
  \colorbox{inkblue}{\makebox[1.5em][c]{\vphantom{Ag}\thepage}}}
\fancypagestyle{plain}{\fancyhf{}\renewcommand{\headrulewidth}{0pt}%
  \fancyfoot[C]{\small\sffamily\bfseries\color{white}%
    \colorbox{inkblue}{\makebox[1.5em][c]{\vphantom{Ag}\thepage}}}}

% caption: nhãn đậm màu accent
\captionsetup{labelfont={bf,sf,color=accent},textfont={small}}

% ---------- BA loại hộp, không thêm loại thứ tư ----------
\tcbset{hopbase/.style={
  enhanced, breakable, sharp corners=downhill, arc=3pt,
  boxrule=0pt, leftrule=3pt, toprule=0pt, bottomrule=0pt, rightrule=0pt,
  left=1.0em, right=1.0em, top=0.65em, bottom=0.65em,
  before skip=1.1em, after skip=1.1em,
  fonttitle=\sffamily\bfseries\small, coltitle=white,
  attach boxed title to top left={xshift=0.9em, yshift=-\tcboxedtitleheight/2},
  boxed title style={sharp corners=downhill, arc=2pt, boxrule=0pt, left=0.55em, right=0.55em,
                     top=0.15em, bottom=0.15em},
  top=1.15em}}
\newtcolorbox{trucgiac}{hopbase, colback=violetbg, colframe=violetdark,
  colbacktitle=violetdark, title={TRỰC GIÁC}, before upper={\bmblock{Trực giác}}}
\newtcolorbox{cambay}{hopbase, colback=redbg, colframe=reddark,
  colbacktitle=reddark, title={CẠM BẪY}}
\newtcolorbox{cannho}{hopbase, colback=bluebg, colframe=steel,
  colbacktitle=steel, title={CẦN NHỚ}, before upper={\bmblock{Cần nhớ}}}

% ---------- tóm tắt đầu tài liệu ----------
% Đặt ngay sau tiêu đề và sau \dependency. 3-6 câu: tài liệu này về gì, đọc xong làm được gì,
% và điều quan trọng nhất nếu chỉ đọc mỗi khối này. KHÔNG phải hộp callout thứ tư.
\newtcolorbox{tomtat}{
  enhanced, breakable, sharp corners=downhill, arc=3pt,
  colback=panelbg, colframe=inkblue,
  boxrule=0pt, leftrule=3pt, toprule=0pt, bottomrule=0pt, rightrule=0pt,
  left=1.0em, right=1.0em, top=1.15em, bottom=0.7em,
  before skip=0.9em, after skip=1.0em,
  fontupper=\small,
  fonttitle=\sffamily\bfseries\small, coltitle=white, colbacktitle=inkblue,
  attach boxed title to top left={xshift=0.9em, yshift=-\tcboxedtitleheight/2},
  boxed title style={sharp corners=downhill, arc=2pt, boxrule=0pt, left=0.55em, right=0.55em,
                     top=0.15em, bottom=0.15em},
  title={TÓM TẮT}, before upper={\bmblock{Tóm tắt}}}

% ---------- mạch vấn đề -> giải pháp ----------
% Mỗi bước là một khối có gạch dọc bên trái, số thứ tự nằm TRONG cột nhãn cùng dòng "Vấn đề",
% các thành phần còn lại (Giải pháp, Giả định, Giá phải trả, Vấn đề mới, Hệ quả) mỗi cái một dòng,
% nhãn thẳng cột, thân thẳng hàng kể cả khi gấp dòng.
\newcounter{machctr}
\newlength{\machlabw}\setlength{\machlabw}{6.6em}
\newlength{\machbase}
\newenvironment{mach}[1][Mạch vấn đề \(\rightarrow\) giải pháp]
  {\par\medskip\noindent
   {\raisebox{0.06em}{\color{accent}\rule{2.6pt}{0.82em}}\hspace{0.5em}%
    \sffamily\bfseries\color{inkblue}#1}\par\nopagebreak\vspace{0.35em}
   \setcounter{machctr}{0}%
   \begin{list}{}%
     {\setlength{\leftmargin}{1.0em}\setlength{\labelwidth}{0pt}\setlength{\labelsep}{0pt}%
      \setlength{\itemsep}{0.95em}\setlength{\topsep}{0.2em}\setlength{\parsep}{0pt}%
      \setlength{\listparindent}{0pt}}%
   \setlength{\machbase}{\leftskip}}
  {\end{list}\medskip}
% một dòng của mạch: #1 = màu nhãn, #2 = nhãn, #3 = 1 nếu ngắt dòng trước
\newcommand{\machrow}[3]{%
  \ifnum#3=1\par\vspace{0.14em}\fi
  \leftskip=\machbase \advance\leftskip by \machlabw
  \noindent\hspace*{-\machlabw}%
  \makebox[\machlabw][l]{\footnotesize\sffamily\bfseries\color{#1}#2}\ignorespaces}
% badge số thứ tự + nhãn "Vấn đề" trên cùng một dòng, không để số lạc giữa hai cột
\newcommand{\machnum}{%
  \tikz[baseline=(n.base)]{\node[circle,fill=inkblue,text=white,inner sep=0pt,minimum size=1.15em,
    font=\scriptsize\sffamily\bfseries] (n) {\arabic{machctr}};}}
\newcommand{\vd}[1]{\stepcounter{machctr}%
  \par\vspace{0.1em}%
  \leftskip=\machbase \advance\leftskip by \machlabw
  \noindent\hspace*{-\machlabw}%
  \makebox[\machlabw][l]{\machnum\hspace{0.4em}\footnotesize\sffamily\bfseries\color{reddark}Vấn đề}%
  \ignorespaces#1\unskip}
\newcommand{\gp}{\machrow{steel}{Giải pháp}{1}}
\newcommand{\gd}{\machrow{tiercol}{Giả định}{1}}
\newcommand{\gia}{\machrow{reddark}{Giá phải trả}{1}}
\newcommand{\vdm}{\machrow{reddark}{Vấn đề mới}{1}}
\newcommand{\hq}{\machrow{tiercol}{Hệ quả}{1}}
\newcommand{\ra}{\ \(\rightarrow\)\ }

% ---------- kiến thức nền cần có trước ----------
% Đặt ngay sau khối tomtat. Định nghĩa NGẮN các khái niệm sơ cấp mà phần thân sẽ dùng,
% để người đọc chưa có nền vẫn theo được. Không phải danh sách link — phải có định nghĩa thật.
% Dùng: \begin{nentang} \kn{learning rate}{độ dài bước ...} \kn{...}{...} \end{nentang}
\newtcolorbox{nentang}{
  enhanced, breakable, sharp corners=downhill, arc=3pt,
  colback=amberbg, colframe=accent,
  boxrule=0pt, leftrule=3pt, toprule=0pt, bottomrule=0pt, rightrule=0pt,
  left=1.0em, right=1.0em, top=1.15em, bottom=0.55em,
  before skip=0.9em, after skip=1.1em,
  fontupper=\small,
  fonttitle=\sffamily\bfseries\small, coltitle=white, colbacktitle=accent,
  attach boxed title to top left={xshift=0.9em, yshift=-\tcboxedtitleheight/2},
  boxed title style={sharp corners=downhill, arc=2pt, boxrule=0pt, left=0.55em, right=0.55em,
                     top=0.15em, bottom=0.15em},
  title={KIẾN THỨC CẦN CÓ TRƯỚC}, before upper={\bmblock{Kiến thức cần có trước}}}
% \kn{thuật ngữ}{định nghĩa một câu}
\newcommand{\kn}[2]{\par\vspace{0.24em}\hangindent=1.15em\hangafter=1\noindent
  {\color{accent}\rule[0.22em]{0.42em}{0.42em}}\hspace{0.45em}\textbf{#1} --- #2\par}

% ---------- macro nội dung ----------
\newcommand{\term}[1]{\textbf{#1}}                 % thuật ngữ lần đầu (đã có tiếng Anh kèm)
\newcommand{\eng}[1]{\emph{#1}}                    % tên tiếng Anh đứng một mình
% \vn{tiếng Việt}{English} — DẠNG CHUẨN cho thuật ngữ ở lần xuất hiện ĐẦU TIÊN trong mỗi file:
%   \vn{hàm mất mát}{loss function}  ->  **hàm mất mát** (*loss function*)
\newcommand{\vn}[2]{\textbf{#1} (\emph{#2})}
% \en{English} — thuật ngữ không nên dịch, giữ nguyên tiếng Anh và in đậm ở lần đầu
\newcommand{\en}[1]{\textbf{\emph{#1}}}
\newcommand{\code}[1]{\texttt{#1}}
% \repo{owner/name} -> link GitHub; \repo{tên trần} -> tra bảng dưới; không có thì để chữ thường.
% \repo[url]{chữ hiện} để chỉ định URL tuỳ ý.
\usepackage{xstring}
% --- hệ sinh thái fiducial marker / hiệu chuẩn / visual servoing ---
\expandafter\def\csname repourl@AprilTag-imgs\endcsname{https://github.com/AprilRobotics/apriltag-imgs}
\expandafter\def\csname repourl@CCTag\endcsname{https://github.com/alicevision/CCTag}
\expandafter\def\csname repourl@COLMAP\endcsname{https://github.com/colmap/colmap}
\expandafter\def\csname repourl@Ceres\endcsname{https://github.com/ceres-solver/ceres-solver}
\expandafter\def\csname repourl@ChromaTag\endcsname{https://github.com/CogChameleon/ChromaTag}
\expandafter\def\csname repourl@DeepArUco\endcsname{https://github.com/AVAuco/deeparuco}
\expandafter\def\csname repourl@DeepTag\endcsname{https://github.com/herohuyongtao/deeptag-pytorch}
\expandafter\def\csname repourl@GTSAM\endcsname{https://github.com/borglab/gtsam}
\expandafter\def\csname repourl@Kalibr\endcsname{https://github.com/ethz-asl/kalibr}
\expandafter\def\csname repourl@OpenCV\endcsname{https://github.com/opencv/opencv}
\expandafter\def\csname repourl@STag\endcsname{https://github.com/bbenligiray/stag}
\expandafter\def\csname repourl@Sophus\endcsname{https://github.com/strasdat/Sophus}
\expandafter\def\csname repourl@TopoTag\endcsname{https://github.com/herohuyongtao/topotag}
\expandafter\def\csname repourl@ViSP\endcsname{https://github.com/lagadic/visp}
\expandafter\def\csname repourl@WhyCon\endcsname{https://github.com/gestom/whycon}
\expandafter\def\csname repourl@apriltag\endcsname{https://github.com/AprilRobotics/apriltag}
\expandafter\def\csname repourl@apriltag-generation\endcsname{https://github.com/AprilRobotics/apriltag-generation}
\expandafter\def\csname repourl@apriltag_ros\endcsname{https://github.com/AprilRobotics/apriltag_ros}
\expandafter\def\csname repourl@aruco_ros\endcsname{https://github.com/pal-robotics/aruco_ros}
\expandafter\def\csname repourl@deltille\endcsname{https://github.com/facebookincubator/deltille}
\expandafter\def\csname repourl@easy_handeye\endcsname{https://github.com/IFL-CAMP/easy_handeye}
\expandafter\def\csname repourl@evo\endcsname{https://github.com/MichaelGrupp/evo}
\expandafter\def\csname repourl@fiducials\endcsname{https://github.com/UbiquityRobotics/fiducials}
\expandafter\def\csname repourl@manif\endcsname{https://github.com/artivis/manif}
\expandafter\def\csname repourl@tagslam\endcsname{https://github.com/berndpfrommer/tagslam}
\makeatletter
\newcommand{\repo}[2][]{%
  \IfSubStr{#2}{/}%
    {\href{https://github.com/#2}{\texttt{\color{steel}#2}}}%
    {\ifx\relax#1\relax
       \StrSubstitute{#2}{\_}{_}[\reponame]%
       \@ifundefined{repourl@\reponame}%
         {\texttt{#2}}%
         {\href{\csname repourl@\reponame\endcsname}{\texttt{\color{steel}#2}}}%
     \else \href{#1}{\texttt{\color{steel}#2}}\fi}}
\makeatother
% nhãn tầng bán rã T1..T4
\newcommand{\T}[1]{{\footnotesize\sffamily\color{tiercol}[T#1]}}
% dòng phụ thuộc ở đầu mỗi tài liệu
\newcommand{\dependency}[1]{%
  \par\vspace{0.15em}\noindent
  \hangindent=1.0em\hangafter=1
  {\footnotesize\sffamily\bfseries\color{white}\colorbox{tiercol}{\vphantom{Ag}PHỤ THUỘC}}\hspace{0.5em}%
  {\footnotesize\sffamily\color{tiercol}#1}\par\vspace{0.7em}}
% câu hỏi tự kiểm ở cuối mỗi tài liệu
\newenvironment{tukiem}{\par\medskip\bmblock{Câu hỏi tự kiểm}%
  {\raisebox{0.06em}{\color{accent}\rule{2.6pt}{0.82em}}\hspace{0.5em}%
   \sffamily\bfseries\color{inkblue}Câu hỏi tự kiểm}%
  {\sffamily\small\color{tiercol}\ \ (trả lời không nhìn bài)}%
  \begin{enumerate}[label=\arabic*.,leftmargin=1.7em]}{\end{enumerate}}

% ---------- hệ thống hoá & ôn tập ----------
% Mức quan trọng: \muc{3} = phải thuộc, \muc{2} = nên nắm, \muc{1} = biết là có, tra khi cần.
\newcommand{\muc}[1]{{\color{accent}\ifcase#1\or\(\bullet\)\or\(\bullet\bullet\)\or\(\bullet\bullet\bullet\)\fi}}

% ACTIVE RECALL cuối mỗi chương — thay cho sơ đồ ôn tập và mạng liên kết khái niệm.
% Hai cột: trái là thuật ngữ / khái niệm / câu hỏi, phải là câu trả lời.
% Che cột phải, tự trả lời, rồi mở ra đối chiếu. r[mức]{hỏi}{đáp}, mức 1..3, mặc định 2.
\newenvironment{onbai}[1][Active recall]
 {\par\medskip\bmblock{Active recall}%
  \noindent{\raisebox{0.06em}{\color{accent}\rule{2.6pt}{0.82em}}\hspace{0.5em}%
   \sffamily\bfseries\color{inkblue}#1}%
  {\sffamily\small\color{tiercol}\ \ (che cột phải, tự trả lời, rồi mở ra đối chiếu)}\par
  \vspace{0.5em}\footnotesize
  \rowcolors{2}{white}{tblalt}%
  \longtable{@{}p{1.75em}L{4.9cm}!{\color{rulecol}\vrule width 0.8pt}L{8.8cm}@{}}
  \toprule
  & \textbf{Thuật ngữ / câu hỏi} & \textbf{Trả lời}\\
  \midrule
  \endfirsthead
  \toprule
  & \textbf{Thuật ngữ / câu hỏi} & \textbf{Trả lời (tiếp)}\\
  \midrule
  \endhead}
 {\bottomrule\endlongtable\par\vspace{0.25em}
  \noindent{\footnotesize\sffamily\color{tiercol}%
  {\color{accent}\(\bullet\bullet\bullet\)} phải thuộc \quad
  {\color{accent}\(\bullet\bullet\)} nên nắm \quad
  {\color{accent}\(\bullet\)} biết là có}\par\medskip}
\newcommand{\ar}[3][2]{{\color{accent}\scriptsize\ifcase#1\or\(\bullet\)\or\(\bullet\bullet\)\or\(\bullet\bullet\bullet\)\fi} & \textbf{#2} & #3\\}
% giữ \y để file chưa kịp chuyển vẫn build được; nó chỉ in một dòng gọn thay vì hàng bảng
\newcommand{\y}[4]{\par\noindent\hangindent=1.2em\hangafter=1
  {\color{accent}\rule[0.22em]{0.42em}{0.42em}}\ \textbf{#2} --- #3\ \ {\color{reddark}\emph{Hỏng khi:}} #4\par}

% ---------- câu hỏi mở đầu & lời giải cuối chương ----------
% khoigoi: đặt ở ĐẦU file chương, sau nentang. 3-5 câu hỏi khiêu khích, mỗi câu phải được
% trả lời trong chương và nhắc lại ở khối traloi cuối chương. Đánh số để đối chiếu.
\newcounter{hoictr}
\newtcolorbox{khoigoi}[1][CÂU HỎI MỞ ĐẦU]{
  enhanced, breakable, sharp corners=downhill, arc=3pt,
  colback=violetbg, colframe=violetdark,
  boxrule=0pt, leftrule=3pt, toprule=0pt, bottomrule=0pt, rightrule=0pt,
  left=1.0em, right=1.0em, top=1.15em, bottom=0.6em,
  before skip=1.0em, after skip=1.1em, fontupper=\small,
  fonttitle=\sffamily\bfseries\small, coltitle=white, colbacktitle=violetdark,
  attach boxed title to top left={xshift=0.9em, yshift=-\tcboxedtitleheight/2},
  boxed title style={sharp corners=downhill, arc=2pt, boxrule=0pt, left=0.55em, right=0.55em,
                     top=0.15em, bottom=0.15em},
  title={#1}, before upper={\setcounter{hoictr}{0}\bmblock{#1}}}
\newcommand{\h}[1]{\par\vspace{0.22em}\stepcounter{hoictr}%
  \hangindent=1.7em\hangafter=1\noindent
  {\sffamily\bfseries\color{violetdark}Q\thehoictr.}\hspace{0.5em}#1\par}
\newtcolorbox{traloi}[1][TRẢ LỜI CÂU HỎI MỞ ĐẦU]{
  enhanced, breakable, sharp corners=downhill, arc=3pt,
  colback=violetbg, colframe=violetdark,
  boxrule=0pt, leftrule=3pt, toprule=0pt, bottomrule=0pt, rightrule=0pt,
  left=1.0em, right=1.0em, top=1.15em, bottom=0.6em,
  before skip=1.0em, after skip=1.1em, fontupper=\small,
  fonttitle=\sffamily\bfseries\small, coltitle=white, colbacktitle=violetdark,
  attach boxed title to top left={xshift=0.9em, yshift=-\tcboxedtitleheight/2},
  boxed title style={sharp corners=downhill, arc=2pt, boxrule=0pt, left=0.55em, right=0.55em,
                     top=0.15em, bottom=0.15em},
  title={#1}, before upper={\setcounter{hoictr}{0}\bmblock{#1}}}
\newcommand{\tl}[1]{\par\vspace{0.22em}\stepcounter{hoictr}%
  \hangindent=1.7em\hangafter=1\noindent
  {\sffamily\bfseries\color{violetdark}A\thehoictr.}\hspace{0.5em}#1\par}

% ---------- kiến thức này đang sống ở đâu (ứng dụng trong công trình SOTA) ----------
\newenvironment{ungdung}[1][Kiến thức này đang sống ở đâu]
 {\par\medskip\bmblock{#1}%
  \noindent{\raisebox{0.06em}{\color{accent}\rule{2.6pt}{0.82em}}\hspace{0.5em}%
   \sffamily\bfseries\color{inkblue}#1}%
  {\sffamily\small\color{tiercol}\ \ (công trình đang dùng chính ý tưởng của chương; chốt 9/2026)}\par
  \vspace{0.45em}\footnotesize
  \rowcolors{2}{white}{tblalt}%
  \longtable{@{}L{3.6cm} L{6.4cm} L{4.6cm}@{}}
  \toprule
  \textbf{Công trình / hệ thống} & \textbf{Dùng ý tưởng nào của chương} & \textbf{Ghi chú}\\
  \midrule
  \endfirsthead
  \toprule
  \textbf{Công trình / hệ thống} & \textbf{Dùng ý tưởng nào của chương} & \textbf{Ghi chú}\\
  \midrule
  \endhead}
 {\bottomrule\endlongtable\par\medskip}
\newcommand{\ud}[3]{#1 & #2 & #3\\}

% ---------- nói lại thật đơn giản (kỹ thuật Feynman) ----------
% Đặt CUỐI mỗi chương, sau traloi. Giải thích lại toàn chương bằng ngôn ngữ thường,
% như đang nói với người chưa học gì, cho tới khi không còn chỗ nào phải nói "cái này phức tạp lắm".
\newtcolorbox{dongian}[1][NÓI LẠI THẬT ĐƠN GIẢN]{
  enhanced, breakable, sharp corners=downhill, arc=3pt,
  colback=greenbg, colframe=greendark,
  boxrule=0pt, leftrule=3pt, toprule=0pt, bottomrule=0pt, rightrule=0pt,
  left=1.0em, right=1.0em, top=1.15em, bottom=0.65em,
  before skip=1.1em, after skip=1.1em, fontupper=\small,
  fonttitle=\sffamily\bfseries\small, coltitle=white, colbacktitle=greendark,
  attach boxed title to top left={xshift=0.9em, yshift=-\tcboxedtitleheight/2},
  boxed title style={sharp corners=downhill, arc=2pt, boxrule=0pt, left=0.55em, right=0.55em,
                     top=0.15em, bottom=0.15em},
  title={#1}, before upper={\bmblock{Nói lại thật đơn giản}}}
% \motcau{...} — câu chốt một dòng, đặt cuối khối dongian
\newcommand{\motcau}[1]{\par\vspace{0.4em}\noindent
  {\sffamily\bfseries\footnotesize\color{greendark}Một câu:}\ \emph{#1}\par}
% \chuadongian{...} — chỗ thật sự không đơn giản hoá được, nói thẳng thay vì giả vờ
\newcommand{\chuadongian}[1]{\par\vspace{0.3em}\noindent
  {\sffamily\bfseries\footnotesize\color{reddark}Chỗ không rút gọn được:}\ #1\par}

% ---------- mini project cuối chương ----------
% Mỗi chương kết thúc bằng một mini project chạy được (mô phỏng hoặc trên bàn thật), nối chuỗi qua cả cây.
% Dùng: \begin{duan}{Tên dự án}{thời lượng}  \dmt{...} \ddl{...} \dbc{...} \dxk{...} \dnv{...}  \end{duan}
\newtcolorbox{duan}[2]{
  enhanced, breakable, sharp corners=downhill, arc=3pt,
  colback=amberbg, colframe=accent,
  boxrule=0pt, leftrule=3pt, toprule=0pt, bottomrule=0pt, rightrule=0pt,
  left=1.0em, right=1.0em, top=1.15em, bottom=0.6em,
  before skip=1.1em, after skip=1.1em, fontupper=\small,
  fonttitle=\sffamily\bfseries\small, coltitle=white, colbacktitle=accent,
  attach boxed title to top left={xshift=0.9em, yshift=-\tcboxedtitleheight/2},
  boxed title style={sharp corners=downhill, arc=2pt, boxrule=0pt, left=0.55em, right=0.55em,
                     top=0.15em, bottom=0.15em},
  title={MINI PROJECT\quad\textperiodcentered\quad #1\quad\textperiodcentered\quad #2},
  before upper={\bmblock{Mini project}\setlength{\machbase}{\leftskip}}}
\newcommand{\duanrow}[2]{%
  \par\vspace{0.16em}%
  \leftskip=\machbase \advance\leftskip by 5.4em
  \noindent\hspace*{-5.4em}%
  \makebox[5.4em][l]{\footnotesize\sffamily\bfseries\color{#1}#2}\ignorespaces}
\newcommand{\dmt}{\duanrow{inkblue}{Mục tiêu}}
\newcommand{\ddl}{\duanrow{tiercol}{Dữ liệu}}
\newcommand{\dbc}{\duanrow{steel}{Các bước}}
\newcommand{\dxk}{\duanrow{reddark}{Xong khi}}
\newcommand{\dnv}{\duanrow{tiercol}{Nối với}}

% ---------- đường dẫn chéo bấm được ----------
% \ptr{B/03-giai-phau-he-thong} -> hiện như mã, bấm vào nhảy tới chương/mục đó.
% Đích được gieo tự động bởi scratchpad/link_paths.py dưới dạng \ptrtarget{<alias>}.
\makeatletter
% \ptrtarget đặt ngay sau \chapter/\section nên \label bắt đúng neo của tiêu đề;
% nhờ vậy cửa sổ xem trước của trình đọc PDF hiện cả tên chương, không bị cắt mất.
\newcommand{\ptrtarget}[1]{\label{tp:#1}}
\newcommand{\ptr}[1]{%
  \@ifundefined{r@tp:#1}{\code{#1}}{\hyperref[tp:#1]{\code{\color{steel}#1}}}}
\makeatother

% ---------- liên kết khái niệm ----------
% \lk{đích}{quan hệ}{giải thích} — liên kết TẠI CHỖ, đặt ngay chỗ khái niệm được nhắc.
%   đích: đường dẫn dạng C/12/03 hoặc B/09.  quan hệ: một trong sáu nhãn dưới đây.
\newcommand{\lk}[3]{%
  \par\vspace{0.2em}\noindent\hangindent=1.6em\hangafter=1
  {\footnotesize\sffamily\bfseries\color{steel}\(\hookrightarrow\)\ #2}%
  {\footnotesize\sffamily\color{tiercol}\ \ \ptr{#1} --- #3}\par\vspace{0.2em}}


% ---------- code ----------
\lstset{
  basicstyle=\ttfamily\footnotesize, breaklines=true, breakatwhitespace=true,
  backgroundcolor=\color{codebg}, frame=leftline, framerule=1.6pt, rulecolor=\color{rulecol},
  xleftmargin=0.9em, framexleftmargin=0.9em, aboveskip=1em, belowskip=1em,
  keywordstyle=\color{inkblue}\bfseries, commentstyle=\color{steel}\itshape,
  stringstyle=\color{reddark}, showstringspaces=false, columns=fullflexible,
  extendedchars=true, literate={á}{{\'a}}1 {à}{{\`a}}1 {ả}{{\h a}}1 {ô}{{\^o}}1 {ế}{{\'\^e}}1
}
\lstdefinestyle{py}{language=Python}

% ---------- thư mục và liên kết ----------
\usepackage[backend=biber,style=ieee,sorting=none,maxbibnames=3,url=false,doi=false,isbn=false]{biblatex}
\usepackage[unicode,colorlinks=true,linkcolor=steel,citecolor=steel,urlcolor=steel,
  bookmarks=true,bookmarksnumbered=true,bookmarksopen=true,bookmarksopenlevel=1,
  bookmarksdepth=4]{hyperref}
\usepackage{bookmark}
% Bookmark đi sâu hơn mục lục in ra: mục lục chỉ tới \section cho gọn,
% còn cây bookmark hiện cả \subsection (các tiểu mục A., B., C.) để nhảy nhanh.
\setcounter{tocdepth}{1}
\bookmarksetup{depth=4, numbered=true, open=true, openlevel=1}

% ---------- bản địa hoá thư mục tài liệu ----------
% biblatex không có vietnamese.lbx nên nó in ra chính tên chuỗi (andothers, jourvol).
% Ánh xạ sang english.lbx rồi ghi đè những chuỗi thật sự xuất hiện sang tiếng Việt.
\DeclareLanguageMapping{vietnamese}{english}
\DefineBibliographyStrings{english}{%
  andothers   = {và cộng sự},
  andmore     = {và nhiều tác giả khác},
  in          = {trong},
  jourvol     = {tập},
  number      = {số},
  volume      = {tập},
  edition     = {ấn bản},
  page        = {tr\adddot},
  pages       = {tr\adddot},
  editor      = {chủ biên},
  editors     = {chủ biên},
  techreport  = {báo cáo kỹ thuật},
  phdthesis   = {luận án tiến sĩ},
  mathesis    = {luận văn thạc sĩ},
  urlseen     = {truy cập},
  version     = {phiên bản},
  july        = {tháng 7},
}

% ---------- bookmark cho các khối cấu trúc ----------
% Mỗi khối (Tóm tắt, Kiến thức cần có trước, Câu hỏi mở đầu, Bảng ôn tập, ...) tự sinh một
% mục bookmark riêng, để nhảy thẳng tới đúng phần cần mà không cuộn trang.
\newcounter{bmblk}
\newcommand{\bmblock}[1]{%
  \stepcounter{bmblk}%
  \bookmark[level=2, dest=bmblk.\thebmblk]{#1}%
  \hypertarget{bmblk.\thebmblk}{}}
\addto\captionsvietnamese{%
  \renewcommand{\partname}{Phần}\renewcommand{\chaptername}{Chương}%
  \renewcommand{\contentsname}{MỤC LỤC}\renewcommand{\tablename}{Bảng}\renewcommand{\figurename}{Hình}}
